% Partie : théorie de la démonstration
% Chapitre : calcul des séquents
% Section : règles et preuves

\documentclass[../../../include/open-logic-section]{subfiles}

\begin{document}

\olfileid[fr]{pt}{seq}{rp}

\olsection{Règles et \usetoken{p}{proof}}

Commençons par quelques définitions et conventions terminologiques.

Considérons, par exemple, les deux règles de \Log{G3c}
qui permettent de prouver des séquents contenant $!A \land !B$
à partir de séquents contenant $!A$ et~$!B$ :
\[
\Axiom$ !A, !B, \Gamma \fCenter \Delta$
\RightLabel{\LeftR{\land}}
\UnaryInf$ !A \land !B, \Gamma \fCenter \Delta$
\DisplayProof
\qquad
\Axiom$\Gamma \fCenter \Delta, !A$
\Axiom$ \Gamma \fCenter \Delta, !B$
\RightLabel{\RightR{\land}}
\BinaryInf$ \Gamma \fCenter \Delta, !A \land !B $
\DisplayProof
\]
Les séquents au-dessus du trait sont les \emph{prémisses},
et celui au-dessous est la \emph{conclusion}.
Les occurrences de formules dans $\Gamma$ et $\Delta$
constituent le \emph{contexte} de chaque règle.
Ce sont les \emph{formules latérales}.
Dans ces schémas, l'occurrence qui s'ajoute au contexte
dans la conclusion est celle de la \emph{formule principale}.
Les occurrences qui s'ajoutent au contexte dans les prémisses
sont celles des \emph{formules actives},
ou \emph{auxiliaires}.

Les règles des quantificateurs sont un peu plus complexes.
Ainsi, les règles de \Log{G1c} pour $\lforall$ sont :
\[\Axiom$ !A(t), \Gamma \fCenter \Delta$
\RightLabel{\LeftR{\lforall}}
\UnaryInf$ \lforall[x][!A(x)],\Gamma \fCenter \Delta$
\DisplayProof
\qquad
\Axiom$ \Gamma \fCenter \Delta, !A(a) $
\RightLabel{\RightR{\lforall}}
\UnaryInf$ \Gamma \fCenter \Delta, \lforall[x][!A(x)]$
\DisplayProof
\]
Dans la règle \LeftR{\lforall}, la formule active est~$!A(t)$,
où $t$ est un terme clos, c'est-à-dire sans variable.
Une telle inférence est conforme au schéma \LeftR{\lforall}
s'il existe une formule~$!A$, une variable~$x$
et un terme clos~$t$ tels que la formule principale
de la conclusion soit $\lforall[x][!A]$ et que la
formule active de la prémisse soit~$\Subst{!A}{t}{x}$.
La règle \RightR{\lforall} se présente de la même façon,
mais, au lieu d'un terme quelconque~$t$, elle exige
dans la prémisse la formule active~$\Subst{!A}{a}{x}$,
où $a$ est une constante. On appelle cette constante
la \emph{constante propre}, ou \emph{eigenvariable},
de l'inférence \RightR{\lforall}.
L'inférence n'est conforme à la règle que si $a$
n'apparaît pas dans le séquent inférieur, donc ni
dans $\Gamma$, ni dans $\Delta$, ni dans~$!A$.
C'est la \emph{condition sur la constante propre}.

Un calcul des séquents tel que \Log{G1c} ou \Log{G3c}
comprend des règles ainsi spécifiées par des schémas,
ainsi que des schémas de séquents qui constituent
ses axiomes. Les preuves de ce calcul sont des
arbres finis de séquents construits inductivement
à partir des axiomes au moyen des règles d'inférence.
Chaque feuille est un axiome, et tout autre séquent
découle des séquents immédiatement au-dessus de lui
par l'une des règles du calcul.

\begin{defn}\ollabel{defn:proof}
Les \emph{preuves} d'un calcul des séquents sont
les arbres finis de séquents définis inductivement
comme suit :
\begin{enumerate}
  \item\ollabel{defn:prf:axiom} Tout axiome est une preuve.
  \item\ollabel{defn:prf:one-prem} Si $\pi_1$ est une preuve
  dont le séquent final est~$S_1$, et si le séquent~$S$
  découle de~$S_1$ par une règle~$R$ à une prémisse, alors
  \[
  \AxiomC{}
  \RightLabel{$\pi_1$}
  \DeduceC{$S_1$}
  \RightLabel{$R$}
  \UnaryInfC{$S$}
  \DisplayProof
  \]
  est une preuve.
  \item\ollabel{defn:prf:two-prem} Si $\pi_1$ et $\pi_2$
  sont des preuves dont les séquents finaux sont~$S_1$
  et~$S_2$, et si le séquent~$S$ découle de~$S_1$
  et~$S_2$ par une règle~$R$ à deux prémisses, alors
  \[
  \AxiomC{}
  \RightLabel{$\pi_1$}
  \DeduceC{$S_1$}
  \AxiomC{}
  \RightLabel{$\pi_2$}
  \DeduceC{$S_2$}
  \RightLabel{$R$}
  \BinaryInfC{$S$}
  \DisplayProof
  \]
  est une preuve.
\end{enumerate}
\end{defn}

\Olref{defn:proof} définit les preuves comme des
arbres de séquents. Nous représentons ces arbres
avec leurs branches dirigées vers le haut.
Les nœuds supérieurs, ou feuilles, sont des axiomes ;
on les appelle les \emph{séquents initiaux}.
Le séquent inférieur est le \emph{séquent final}.
Si $\pi$ est une preuve de séquent final~$S$,
on écrit aussi $\pi \Proves S$.
S'il existe une preuve du séquent~$S$, on écrit
$\Proves S$, en précisant éventuellement le calcul
à gauche du symbole~$\Proves$ : par exemple,
$\Log{G1c} \Proves !A, !B \Sequent !A \land !B$.
Tout séquent d'une preuve, sauf les séquents initiaux,
est la \emph{conclusion} d'une inférence obtenue par
une règle~$R$. Les \emph{prémisses} de cette inférence
sont les séquents immédiatement au-dessus
de sa conclusion.

Les preuves utilisées dans la construction inductive
d'une preuve en sont les \emph{sous-preuves}.
Celles dont les séquents finaux sont les prémisses
d'une inférence sont les \emph{sous-preuves immédiates}
de cette inférence.

Il faut souvent mesurer la complexité d'une preuve
par un entier naturel. On pourrait compter ses séquents,
mais une mesure souvent plus utile est sa hauteur :
le nombre maximal d'inférences sur un chemin reliant
le séquent final à un séquent initial. On compte donc
les arêtes de ce chemin ; une preuve réduite à un axiome
a hauteur zéro. Cette notion se définit aussi inductivement.

\begin{defn}
La \emph{hauteur}~$\pheight{\pi}$ d'une preuve~$\pi$
est définie inductivement comme suit :
\begin{enumerate}
  \item Si $\pi$ se réduit à un axiome, $\pheight{\pi} = 0$.
  \item Si $\pi$ se termine par une inférence à une prémisse,
  de sous-preuve immédiate~$\pi_1$, alors
  $\pheight{\pi} = \pheight{\pi_1} + 1$.
  \item Si $\pi$ se termine par une inférence à deux prémisses,
  de sous-preuves immédiates~$\pi_1$ et~$\pi_2$, alors
  $\pheight{\pi} = \max(\pheight{\pi_1}, \pheight{\pi_2}) + 1$.
\end{enumerate}
Si le séquent~$S$ possède une preuve de hauteur $\le n$,
on écrit $\Proves[n] S$.
\end{defn}

\begin{ex}
Dans \Log{G3c}, tout séquent de la forme
$!C, \Gamma \Sequent \Delta, !C$ est un axiome
si $!C$ est atomique. Supposons que $!D$ et $!E$
soient des formules atomiques closes.
Les séquents $!D, !E \Sequent !D$ et
$!D, !E \Sequent !E$ sont alors des instances
du schéma $!C, \Gamma \Sequent \Delta, !C$.
Par exemple, en prenant $!C \ident !E$,
$\Gamma = \{!D\}$ et $\Delta = \emptyset$
dans $!C, \Gamma \Sequent \Delta, !C$,
on obtient précisément $!D, !E \Sequent !E$.
Rappelons que les contextes sont des multiensembles :
$!D, !E \Sequent !E$ et $!E, !D \Sequent !E$
sont donc le même séquent, puisque seul l'ordre
des formules de l'antécédent a changé.

Le séquent $!D, !E \Sequent !D$, instance d'un axiome
de~\Log{G3c}, constitue à lui seul une preuve
de~\Log{G3c}. Il en va de même de
$!D, !E \Sequent !E$, par la
\olref{defn:proof}\olref{defn:prf:axiom}.
Appelons ces preuves $\pi_1$ et~$\pi_2$.
Par \olref{defn:prf:two-prem}, on peut les combiner
pour construire une nouvelle preuve, $\pi_3$,
au moyen de la règle à deux prémisses~$\RightR{\land}$ :
\[
\Axiom$!D, !E \fCenter !D$
\Axiom$!D, !E \fCenter !E$
\RightLabel{\RightR{\land}}
\BinaryInf$!D, !E  \fCenter !D \land !E $
\DisplayProof
\]
À partir de $\pi_3$, on obtient à son tour
la preuve~$\pi_4$
\[
\Axiom$!D, !E \fCenter !D$
\Axiom$!D, !E \fCenter !E$
\RightLabel{\RightR{\land}}
\insertBetweenHyps{\hspace{5ex}}
\BinaryInf$!D, !E  \fCenter !D \land !E $
\LeftSubproofLabel{$\pi_3$}
\RightLabel{\LeftR{\land}}
\UnaryInf$!D \land !E  \fCenter !D \land !E$
\RightSubproofLabel{$\pi_4$}
\DisplayProof
\]
par la règle à une prémisse~$\LeftR{\land}$,
en appliquant \olref{defn:proof}\olref{defn:prf:one-prem}.
La sous-preuve immédiate de la dernière inférence
de~$\pi_4$ est~$\pi_3$. Les sous-preuves immédiates
de l'inférence $\RightR{\land}$ sont $\pi_1$ et~$\pi_2$.
Les sous-preuves de $\pi_4$ sont $\pi_1$, $\pi_2$,
$\pi_3$ et $\pi_4$ elle-même.
On a $\pheight{\pi_1} = \pheight{\pi_2} = 0$,
$\pheight{\pi_3} = 1$ et $\pheight{\pi_4} = 2$.
On a donc montré que
$\Log{G3c} \Proves[4] !D \land !E  \fCenter !D \land !E$
pour tous $!D$ et~$!E$ atomiques.
La borne quatre est valable, mais non optimale :
la preuve construite a hauteur deux.
\end{ex}

\end{document}
